\renewcommand{\footerfilename}{Apple.tex}

\section{Apple MacOS Mx Processors}

IRAF was developed in the age of command line processing at the time when the
GUI interfaces were starting to emerge. The X11 standard used two programs,
the nomenclature was reversed from today's meanings. The X-server was a 
simple user program operating somewhere on the network -- either within
the localhost or a remote node. The X-Client was a generic program that
understood X-widget and X-gadget\footnote{A gadget handles the drawing
of something where  gadget with a window-manager wrapper is referred to
as a widget} and the content for that widget was prepared by the X-Server,
broken-down into simple commands, and sent across a communication link
to the X-client. The X-client then 'rendered' the result for the user.
This converted massive images of screens down to very small packets.
The X-server and X-client agreed on a widget-set. So I could design a
box around text -- and the X-Server only had to supply the handful of
characters in one packet.  Highly efficent, to the point that complex
CAD operations could be handled over a 1200 baud modem. I did that.
It was a bit painful but dooable. 

Now IRAF comes along and uses X11 to help with this process. At the
time it was made, fonts were few and far between. Terminals were
the rage, and the \dhl{termcap} file entries handled how the program
sent/recieved and rendered results. This was more to help the computer
talk to a terminal that with a direct connection. Similar files existed
for plotters and graphics. Properitary protocols abounded, but played
close enough to the Unix termcap and other -cap files to be workable.
The FITS format was nailed down fairly early and did not have to
waiver much. Astronomers do not care about these details, they want
to load a 'program' and use it.

We have strayed from this ease of installation and use.

The current Linux file system -- applicable to all platforms,
calls for a \dhl{/opt} directory for a 'vendor', under which
that vendor is able to put all of its files/programs. 

To this end IRAF does is a specialized program suite within an
extremely narrow market. It is unworthy of having its nitch in /etc.

/opt/iraf-xxx/etc is a perfectly good place to live.

See \url{https://refspecs.linuxfoundation.org/FHS_3.0/fhs-3.0.pdf} and
\url{https://refspecs.linuxfoundation.org/FHS_3.0/fhs/index.html} for
details of the filesystem.

The \dhl{/usr} directory is essentially reserved for system wide,
common programs. We require Python virtual environments at this time.



This note sequence relates to trying to install a working IRAF/PyRAF onto 
a Macbook 2024 M4 machine under Sequoia OS. 

The community edition does not install and work, and in the process liters
the main system directories with many files. 

The NOIRLabs edition, installs but the epar command does not work.

Numerous additions to IRAF/PyRAF have been made, mainly with an eye to
improved graphics. There are differences with the Python readline command
that is a bit noxious.

Currently various editions pollute 

\subsection{NOIRLabs on Mac M4}

There have been significant changes with the placement of  the basic
user's iraf directory, its contents and questionable changes to login.cl.





In the past, login.cl only handled what the main iraf system needed
to operate -- with flexibility deferred to loginuser.cl.

The M1,M2,M3 and now M4 processors are supported by IRAF Comminuty, but with a 
great deal of difficulties.


\subsection{Apple Oddities}


\subsubsection{BASH}

Use \dhl{type -f <commandname>} to see what kind of file/alias/function
\dhl{commandname} might be,

Use \dhl{echo \$BASH\_VERSINFO} to have the shell self-report what version
it is. This is handy in the event some alias has put you in a weird place.
For example: starting pyraf under /opt/homebrew/bin/bash (version 5.2.37(1)-release
and the using the \dhl{! echo  \$BASH\_VERSINFO} you find yourself back in
the system bash. Where this switch takes place is unknown to me 2025-04-18T12:13:30-0600,
\ltodo{bash}{Figure out the switch in bash versions after starting pyraf.}



The Apple Bash is Version 3.2.57(1)-release, where the current Bash is in the 
5.1.16(1)-release range. Apple has different executables for /bin/sh and /bin/bash
but both seem to be the same version. The /bin/sh is smaller.

The IRAF-Community scripts all refer to /bin/sh
\begingroup \fontsize{10pt}{10pt}
\selectfont
%%\begin{Verbatim} [commandchars=\\\{\}]
\begin{verbatim} 
./unix/hlib/irafcl.sh
./unix/hlib/mkiraf.sh
./unix/hlib/setup.sh
./unix/hlib/mkmlist.sh
./unix/hlib/f77.sh
./unix/hlib/irafarch.sh
./unix/hlib/mkfloat.sh
./unix/f2c/trim_f2c.sh
\end{verbatim}
\endgroup
%% \end{Verbatim}

There are no scripts in pyraf.


